\section{Primary-source corrections for the incoming report}
\label{sec:web-primary-corrections}

This is a claim-specific source audit of the supplied report, with an independent
operator calculation. It does not certify the complete proofs of the cited papers.
The downloaded snapshots and their hashes are listed in
\texttt{retrieval\_receipt.json}; the inspection date is 8 September 2026.

\subsection{The fractional theorem and the exact endpoint statement}
The published Theorem 1 of C\'ordoba, Mart\'inez-Zoroa and Zheng,
\href{https://doi.org/10.1007/s00205-026-02198-0}{ARMA 250 (2026), article 38},
uses $|\nabla|^\alpha=(-\Delta)^{\alpha/2}$ and includes
$0\leq\alpha<(22-8\sqrt7)/9$.
It asserts a solution on $\mathbb R^3\times[0,1]$ with positive viscosity,
forcing in $L^1([0,1];C^{1,\epsilon}_x)\cap L^\infty([0,1];L^2_x)$
for some $\epsilon>0$, velocity in $L^\infty_{t,x}\cap L^\infty_tL^2_x$,
spatial smoothness for $t<1$, and
$\lim_{t\uparrow1}\int_0^t\|\nabla u(\cdot,s)\|_\infty\,ds=\infty$.
Section 1.1 says both solution and forcing are smooth before blow-up;
Section 4.4 gives the corresponding preterminal construction.
The theorem does not assert all-order mixed smoothness through $t=1$.
The incoming report's positive range is valid as a subrange, but omits the
endpoint $\alpha=0$ from the source theorem. Its distinction between
$(-\Delta)^s$ and $|\nabla|^\alpha$ is correct: $s=\alpha/2$.
The cumulative section on the dissipative vortex-layer programme already
contains the parameter calculation and viscosity transfer, so it should be
retained instead of duplicated here.

\subsection{Tao's averaging: exact operator and dynamics}
In \href{https://arxiv.org/html/1402.0290v3}{arXiv:1402.0290v3},
equations (1.11)--(1.13) include dilations, rotations, and real order-zero
Fourier multipliers. Footnote 8 records the role of dilation averaging in the
non-degeneracy argument. The text preceding (1.16) explicitly says that
arbitrary such averaging does not automatically preserve energy cancellation.
Theorem 1.5 chooses a symmetric operator with cancellation and Schwartz
divergence-free data for which no global $H^{10}$ mild solution exists.
Thus the report's broad description is supported, but the dilation and
cancellation qualifications should be supplied. The following calculation
proves the precise relation with classical forced dynamics.

Retain the source's Fourier convention
$\widehat u(\xi)=\int_{\mathbb R^3}u(x)e^{-2\pi i x\cdot\xi}\,dx$.
Let $H=H^{10}_{\mathrm{df}}(\mathbb R^3)$ be the real divergence-free Sobolev
space, with norm
$\|u\|_H^2=\int(1+4\pi^2|\xi|^2)^{10}|\widehat u(\xi)|^2\,d\xi$.
Its $L^2$ dual is denoted $H^*$.
The Leray projection $P$ has symbol
$I-\xi\otimes\xi/|\xi|^2$ for $\xi\ne0$; its value at zero is immaterial.
Define
\[
 B(u,v)=-\frac12 P\big((u\cdot\nabla)v+(v\cdot\nabla)u\big).
\]
For a probability space $(\Omega,\mu)$, retain
$R_{j,\omega}\in SO(3)$, $C^{-1}\leq\lambda_{j,\omega}\leq C$ with $C\geq1$,
and real order-zero symbols $m_{j,\omega}$ satisfying
$m_{j,\omega}(-\xi)=\overline{m_{j,\omega}(\xi)}$ and
\[
 \|m\|_k=\sup_{\xi\ne0}|\xi|^k|\nabla^k m(\xi)|<\infty,
 \qquad \int_\Omega\prod_{j=1}^3\|m_{j,\omega}\|_{k_j}\,d\mu<\infty
 \quad(k_j\geq0).
\]
Write
\[
 \operatorname{Rot}_R u(x)=Ru(R^{-1}x),\qquad
 \operatorname{Dil}_\lambda u(x)=\lambda^{3/2}u(\lambda x),\qquad
 A_{j,\omega}=m_{j,\omega}(D)\operatorname{Rot}_{R_{j,\omega}}
                       \operatorname{Dil}_{\lambda_{j,\omega}}.
\]
Retain the measurability of the source's parameter families: the maps
$\omega\mapsto R_{j,\omega}$ and $\omega\mapsto\lambda_{j,\omega}$
are measurable, and $(\omega,\xi)\mapsto m_{j,\omega}(\xi)$ is jointly
measurable. In particular, for every fixed $u\in H$ the map
$\omega\mapsto A_{j,\omega}u$ is strongly measurable in $H$.
Here is a direct justification of this last assertion. Its Fourier
representative is jointly measurable by the displayed formulas for
rotation, dilation and multiplication. Its scalar coordinates against a
countable orthonormal basis of the weighted Fourier $L^2$ space are
measurable integrals. The finite orthogonal projections onto that basis
are measurable finite-dimensional maps and converge in $H$ at every
$\omega$, proving strong measurability. The same argument in unweighted
$L^2$, using inverse rotations and dilations and conjugate symbols,
proves that $\omega\mapsto A_{j,\omega}^*g$ is strongly measurable
for every fixed $g\in L^2$. These are measurability hypotheses and their
consequences; norm integrability alone is not being used to infer them.
The averaging is the bilinear map $H\times H\to H^*$ specified by
\[
 \langle\widetilde B(u,v),w\rangle
 =\int_\Omega\langle B(A_{1,\omega}u,A_{2,\omega}v),A_{3,\omega}w\rangle
 \,d\mu(\omega).                                      \tag{WA1}
\]
All three operators preserve divergence. For rotations this follows by the
chain rule and $R^{-1}R=I$; for dilations divergence is multiplied by
$\lambda^{5/2}$; scalar Fourier multipliers commute with each derivative.
They preserve real fields by the stated symbol symmetry. Changing variables
$y=\lambda x$, and then $y=R^{-1}x$, proves that dilations and rotations are
$L^2$ isometries with respective adjoints $\operatorname{Dil}_{\lambda^{-1}}$
and $\operatorname{Rot}_{R^{-1}}$. Plancherel gives
$m(D)^*=\overline m(D)$. Consequently
\[
 \widetilde B(u,v)=\int_\Omega A_{3,\omega}^*
                  B(A_{1,\omega}u,A_{2,\omega}v)\,d\mu(\omega),
 \quad
 A_3^*=\operatorname{Dil}_{\lambda_3^{-1}}
       \operatorname{Rot}_{R_3^{-1}}\overline m_3(D).       \tag{WA2}
\]

Here the integral actually defines an $L^2$ field, as follows without an
unstated convergence step. Fourier Cauchy--Schwarz gives
$\|u\|_\infty\leq C_S\|u\|_H$, where
\[
 C_S=\left(\int_{\mathbb R^3}(1+4\pi^2|\xi|^2)^{-10}\,d\xi\right)^{1/2}<\infty;
\]
the radial integrand at infinity is bounded by a constant times $r^{-18}$.
Also $\|\nabla u\|_2\leq\|u\|_H$ and $\|P\|_{2\to2}=1$.
Therefore $\|B(u,v)\|_2\leq C_S\|u\|_H\|v\|_H$.
The Fourier formula for dilation is
$\widehat{\operatorname{Dil}_\lambda u}(\xi)=\lambda^{-3/2}\widehat u(\xi/\lambda)$,
which gives $\|A_j u\|_H\leq C^{10}\|m_j\|_0\|u\|_H$.
The integrand in (WA2) thus has $L^2$ norm at most
$C_SC^{20}\prod_j\|m_j\|_0\|u\|_H\|v\|_H$.
It is also strongly measurable: continuity of $B:H\times H\to L^2$
makes $g_\omega=B(A_{1,\omega}u,A_{2,\omega}v)$ strongly measurable.
Approximate $g_\omega$ pointwise by measurable simple $L^2$-valued
functions. Applying $A_{3,\omega}^*$ to each simple function gives a
strongly measurable function by the preceding fixed-vector result.
The images converge pointwise in $L^2$, since each $A_{3,\omega}^*$
has finite operator norm. Thus the integrand is strongly measurable.
Together with its integrable norm bound this proves existence of its
Bochner integral, the asserted map, and the duality identity (WA1).
No operator is replaced by an equal-looking different operator.

For smooth fields, the Laplacian obeys
\[
 \Delta A_j=\lambda_j^2 A_j\Delta,                       \tag{WA3}
\]
because it commutes with scalar Fourier multipliers and rotations, and each
of the two derivatives of $u(\lambda_jx)$ contributes a factor $\lambda_j$.
Thus the actual dissipative operator remains explicitly present.

Fix the particular cancellation-preserving $\widetilde B$ from Tao's theorem.
For any $u\in C(I;H)\cap C^1(I;L^2)$ on a time interval $I$, set
\[
 f_u=(\widetilde B-B)(u,u),\qquad
 p_u=-\Delta^{-1}\partial_i\partial_j(u_i u_j).
\]
Both fields are well defined in $L^2$ at each time: $u_i u_j\in L^2$ by
$\|u\|_\infty\|u\|_2<\infty$, and the scalar pressure multiplier is bounded.
Set $N=(u\cdot\nabla)u$. Divergence freedom gives
$N_i=\partial_j(u_j u_i)$, hence
$\nabla p_u=-(I-P)N$ in distributions and in $L^2$.
Since $B(u,u)=-PN$, the following is an identity of $L^2$ residuals:
\[
 \partial_tu+N+\nabla p_u-\nu\Delta u-f_u
 =\partial_tu-\nu\Delta u-\widetilde B(u,u),\qquad \nu>0. \tag{WA4}
\]
Indeed, $N+\nabla p_u=PN=-B(u,u)$; substitution proves every sign in
(WA4). In particular the source's coefficient $\nu=1$ is retained when
applying this identity to its solution. The map $u\mapsto(u,p_u,f_u)$
transfers its averaged equation exactly to the ordinary forced equation,
with the same spatial coordinates and the same time variable. Conversely
the forced equation with this explicitly prescribed $f_u$ gives the averaged
equation by the same identity. The force is divergence free because both
bilinear operators take values in divergence-free fields.

There is also an exact energy relation. For $u\in H$, integration against a
cutoff $\chi_R(x)=\chi(x/R)$, equal to one on $|x|\leq R$ and supported on
$|x|\leq2R$, gives
\[
 \int\chi_R N\cdot u
 =-\frac12\int |u|^2u\cdot\nabla\chi_R.
\]
The right side is bounded in absolute value by
$\|\nabla\chi\|_\infty\|u\|_\infty\|u\|_2^2/(2R)$ and tends to zero.
Dominated convergence applies on the left because $N\in L^2$ and $u\in L^2$.
As $P$ is self-adjoint and $Pu=u$, this proves
$\langle B(u,u),u\rangle=0$. The chosen averaging has the same cancellation,
so $\langle f_u,u\rangle=0$. Along a trajectory with zero residual in (WA4),
Plancherel gives $\langle\Delta u,u\rangle=-\|\nabla u\|_2^2$ and therefore
\[
 \frac{d}{dt}\frac12\|u(t)\|_2^2=-\nu\|\nabla u(t)\|_2^2. \tag{WA5}
\]
This proves the force-defect bridge and its energy behavior. It supplies no
all-order terminal estimates for $f_u$; those estimates are not a conclusion
of (WA4) or (WA5). Consequently this exact transfer does not itself establish
the requested smooth-through-terminal-time forced singularity.

\subsection{The inspected numerical Euler manuscript's actual status}
The author-hosted
\href{https://anima-ai.org/wp-content/uploads/2026/09/Euler.pdf}{Ganeshram--Duruisseaux--Anandkumar manuscript}
has 107 pages in the downloaded snapshot. Its abstract describes an
approximate profile and evidence for stability. Page 5 explicitly identifies
unfinished certification of constants and margins, allowing refinement of
estimates. Page 13, Section 4.6, additionally identifies local well-posedness
and continuation in rescaled time as necessary for an unconditional blow-up
result. Page 14, Section 4.9, leaves terms outside the signed matrices to be
certified. Thus the report's candidate-status description is supported;
the remaining work should not be described solely as evaluating constants.
This audit checks these status statements, not their eventual feasibility or
the correctness of all individual certificates.
Page 7, equations (3)--(7), defines $\lambda$ as the common radial and axial
length exponent, and retains a moving axial centre with
$\dot z_c=-C(T-t)^{\lambda-1}$.
The report's fixed-centre illustrative ansatz therefore does not reproduce
all the source's coordinates. The exact length dictionary is
$\beta=\lambda$. Here is the complete coordinate and equation correspondence,
including the axial translation and its velocity.

Let $I\subset\mathbb R$ be an interval, $\nu\geq0$, and
$d\in C^\infty(I;\mathbb R^3)$. Write
$\mathcal N_\nu(u,p)=\partial_tu+(u\cdot\nabla)u+\nabla p-\nu\Delta u$
on $\mathbb R^3\times I$, and let $\mathcal S_\nu(I)$ be the set of smooth
triples $(u,p,f)$ on this domain with $u,f$ vector valued, $p$ scalar valued,
$\operatorname{div}u=0$, and $\mathcal N_\nu(u,p)=f$.
No spatial decay restriction is imposed in this definition. Define
\begin{align}
 u_d(x,t)&=u(x-d(t),t),&p_d(x,t)&=p(x-d(t),t),\notag\\
 f_d(x,t)&=f(x-d(t),t)
       -\sum_{j=1}^3d'_j(t)\partial_j u(x-d(t),t).
 \label{eq:wb-moving-translation}
\end{align}
These formulas define a bijection $\mathcal S_\nu(I)\to\mathcal S_\nu(I)$.
Indeed $\partial_tu_d=(\partial_tu-d'\cdot\nabla u)(x-d,t)$,
$\partial_{x_i}u_d=(\partial_i u)(x-d,t)$, and
$\partial_{x_i}\partial_{x_j}u_d=(\partial_i\partial_j u)(x-d,t)$.
Consequently divergence, convection, pressure gradient, and Laplacian
are the corresponding original quantities evaluated at $(x-d,t)$, while
the displayed time derivative gives precisely $f_d$.
The inverse, including its force component, is
\[
 u(y,t)=u_d(y+d(t),t),\quad p(y,t)=p_d(y+d(t),t),\quad
 f(y,t)=f_d(y+d(t),t)+d'(t)\cdot\nabla u_d(y+d(t),t).
\]
Substitution in both orders proves the identity map on every component.
The spatial Jacobian is the identity matrix with determinant one, so
$\|u_d(t)\|_2=\|u(t)\|_2$ and
$\|\nabla u_d(t)\|_2=\|\nabla u(t)\|_2$ whenever these norms are finite.
Every spatial derivative and every compact support is translated exactly.
Time derivatives retain derivatives of $d$; no bound through an endpoint
is inferred when these derivatives are unbounded there.

For clarity the velocity expressed in the moving frame is a second, also
exact, map. For $(u,p,f)\in\mathcal S_\nu(I)$ put
\[
 v(y,t)=u(y+d(t),t)-d'(t),\quad
 q(y,t)=p(y+d(t),t),\quad g(y,t)=f(y+d(t),t)-d''(t).
\]
Here $\partial_tv=(\partial_tu+d'\cdot\nabla u)(y+d,t)-d''$ and
$(v\cdot\nabla)v=((u-d')\cdot\nabla u)(y+d,t)$; the two $d'$ terms
cancel exactly, proving $\mathcal N_\nu(v,q)=g$ and
$\operatorname{div}v=0$. Its inverse is
$u(x,t)=v(x-d,t)+d'$, $p(x,t)=q(x-d,t)$,
$f(x,t)=g(x-d,t)+d''$.
Equivalently $q_a=q+d''\cdot y$ gives
$\mathcal N_\nu(v,q_a)=f(y+d,t)$, with inverse pressure
$p(x,t)=q_a(x-d,t)-d''\cdot(x-d)$.
The affine pressure and constant velocity are retained, not discarded.
If $u(t)\in L^2(\mathbb R^3)$ and $d'(t)\ne0$, then $v(t)\notin L^2$:
otherwise the difference of two $L^2$ fields would be the nonzero constant
$d'$, whose squared integral on balls is $|d'|^2(4\pi R^3/3)\to\infty$.
Thus the first translation map preserves finite energy; the second has
the explicitly proved different velocity and pressure behavior.

Now retain $T,\tau_*,L_*,U_*>0$, $x_*\in\mathbb R^3$, and
$\beta\in\mathbb R$. For $0\leq t<T$ put
$s=(T-t)/\tau_*$ and $\ell=L_*s^\beta>0$.
For a fixed smooth profile pair $(U,P)$ on $\mathbb R^3$ with
$\operatorname{div}_Y U=0$, set
\[
 Y=\frac{x-x_*-d(t)}{\ell},\qquad
 u_d=U_*s^{\beta-1}U(Y),\qquad
 p_d=U_*^2s^{2\beta-2}P(Y).
\]
At each fixed $t$, $Y\mapsto x=x_*+d(t)+\ell Y$ is a diffeomorphism
$\mathbb R^3\to\mathbb R^3$ with Jacobian $\ell I_3$ and determinant
$\ell^3$. The field inverse on its image is
$U(Y)=U_*^{-1}s^{1-\beta}u_d(x_*+d+\ell Y,t)$ and
$P(Y)=U_*^{-2}s^{2-2\beta}p_d(x_*+d+\ell Y,t)$.
Because $\partial_tY=\beta Y/(\tau_*s)-d'/\ell$,
$\partial_{x_i}Y_j=\delta_{ij}/\ell$, direct differentiation of each
component gives the complete prescribed force
\begin{align}
 f_d={}&\frac{U_*}{\tau_*}s^{\beta-2}
 \left[(1-\beta)U+\beta(Y\cdot\nabla_Y)U
 +\frac{U_*\tau_*}{L_*}
       \{(U\cdot\nabla_Y)U+\nabla_Y P\}
 -\frac{\nu\tau_*}{L_*^2}s^{1-2\beta}\Delta_Y U\right]\notag\\
 &-\frac{U_*}{L_*}s^{-1}(d'(t)\cdot\nabla_Y)U.
 \label{eq:wb-moving-profile-force}
\end{align}
All profiles on the right are evaluated at $Y$. The first line is the
fixed-centre residual evaluated in translated coordinates; the second is
exactly the force correction in \eqref{eq:wb-moving-translation}.
Conversely the displayed field inverse and inverse force translation
recover the fixed-centre forced equation with its original viscosity.

In the cited axial convention put $d(t)=z_c(t)e_3$ and retain the source's
constant $C$ and exponent $\lambda$, so
$z_c'=-C(T-t)^{\lambda-1}$. For every reference $t_0<T$ the full primitive is
\[
 z_c(t)=
 \begin{cases}
 z_c(t_0)+\dfrac C\lambda\big((T-t)^\lambda-(T-t_0)^\lambda\big),
     &\lambda\ne0,\\[3pt]
 z_c(t_0)+C\log\!\dfrac{T-t}{T-t_0},&\lambda=0.
 \end{cases}
\]
Differentiation proves the required negative sign in both cases.
When $\beta=\lambda$, the last line of
\eqref{eq:wb-moving-profile-force} is
\[
 \frac{CU_*\tau_*^{\lambda-1}}{L_*}
 s^{\beta-2}\partial_{Y_3}U
 =\frac{U_*}{\tau_*}s^{\beta-2}
       \frac{C\tau_*^\lambda}{L_*}\partial_{Y_3}U.
\]
This is the exact retained drift in the moving-centre profile equation.
The arbitrary initial centre is present in the inverse coordinate map.
All formulas hold on the preterminal domain $s>0$, including $C=0$;
they make no unproved assertion of an extension through $s=0$.
They relate the two coordinate presentations and their forced equations;
they do not identify their profile solutions or discard the source's
additional profile equations.
